\section{Security and Threat Surface}
\label{sec:security}

This section synthesises the security and threat-modelling
literature on inter-agent protocols and applies the twelve-class
threat catalogue introduced as Axis 8 of the rubric
(\S\ref{sec:rubric:axis8}) protocol-by-protocol.
Section~\ref{sec:security:model} fixes a composite threat model
covering the four principal entities of a contemporary agentic
deployment (user, orchestrator, remote agent, tool or MCP
server) and identifies where on that model each of the twelve
threat classes manifests.
Section~\ref{sec:security:identity} to \ref{sec:security:tee}
each treat one major threat category in depth, drawing on the
specification-level coverage in
Table~\ref{tab:threat-surface} and the empirical and analytical
literature catalogued in our references.
Section~\ref{sec:security:open} closes with the unresolved
problems that should drive the next standardisation cycle.
Half of these reappear in \S\ref{sec:openproblems} as the
top-priority items on the survey's open-problems agenda.

This is deliberately \emph{not} a re-derivation of the security
literature (\S\ref{sec:intro:notwhat}, D3). Kong et
al.~\cite{kong2025communication} already provide a 48-page
treatment that is the canonical primary source for the field.
What \S\ref{sec:security} adds is the spec-coverage view: a
matrix that says, for each protocol's published
specification at the 2026-Q1 cutoff, which threats the spec
addresses with normative language, which it acknowledges only
in non-normative recommendations, and which it leaves out of
scope entirely.

\subsection{Composite threat model}
\label{sec:security:model}

The threat model spans four entities and the three edges
between them. The model is consistent with the layered
architectures of
Kong~\cite{kong2025communication}, A2ASecBench~\cite{a2asecbench},
and the threat-modelling survey~\cite{threatmodel2026}, and
serves as the geometry against which we locate each of the
twelve threat classes.

\begin{description}
  \item[Entity 1: User principal.] A human or
    organisational identity on whose behalf an orchestrator
    is acting. The principal authorises specific intents and
    delegates a bounded mandate to the orchestrator. Threats:
    cross-tenant identity confusion, principal-impersonation
    via stolen mandates, mandate-replay across delegation
    boundaries.
  \item[Entity 2: Orchestrator agent.] The primary
    agent (typically an LLM with an associated runtime)
    that holds the principal's mandate and selects which
    remote agents or tools to invoke. Threats:
    prompt injection at the principal-orchestrator
    boundary, orchestrator output validating
    a downstream attack, side-channel leakage through
    orchestrator logs.
  \item[Entity 3: Remote agent.] A peer agent reached
    via A2A, ANP, AGNTCY/SLIM, Coral, or an equivalent
    peer-messaging substrate. Threats: agent impersonation,
    AgentCard supply-chain attacks, capability escalation
    through ambiguous Skill descriptions, registry poisoning
    against the discovery convention.
  \item[Entity 4: Tool or MCP server.] A
    capability-providing endpoint reached via MCP or an
    equivalent tool-binding protocol. Threats: tool
    poisoning, indirect prompt injection through tool
    results, malicious schema in tool input contracts,
    confused-deputy attacks where the orchestrator's
    OAuth scope is exceeded by tool calls it doesn't
    authorise.
\end{description}

The three edges carry the messages on which all
twelve threats land: principal-to-orchestrator,
orchestrator-to-remote-agent, and
orchestrator-to-tool. Twelve threats is the smallest count
that lets every one of the six published threat-modelling
papers~\cite{kong2025communication, a2asecbench,
threatmodel2026, protocol_exploits, a2a_weaknesses,
zhan2024injecagent} map onto a single matrix without
hiding a threat under an unrelated heading; the granularity
is set by the literature, not by the survey's preferences.

\clearpage

Table~\ref{tab:threat-surface} presents the matrix.
Eight of the twelve threats land on multiple edges
simultaneously. Registry poisoning, for example, is an
attack against the discovery infrastructure. Its impact
flows along every edge that consumes a poisoned descriptor.
The matrix records the protocol's spec-level coverage of
each threat regardless of which edges that protocol mediates,
with the \texttt{na} marker reserved for threats that are
structurally outside the protocol's scope (e.g., MCP has no
AgentCard, so the AgentCard supply-chain row is
\texttt{na} for MCP).

% Generated by scripts/gen_table3.py: do NOT hand-edit.
% Source: data/threat-coverage.yaml.

\begin{table}[h!]
  \centering
  \caption{Threat-surface coverage matrix. Twelve threat classes (rows) $\times$ eight protocols (cols). Cells: \textbf{A} = Addressed at the spec level; P = Partial or non-normative; N = Not addressed; n/a = Not applicable. Every cell is scored against specification text. ACP and AGNTCY were rescored from primary sources in 2026-Q3: ACP against its OpenAPI document, which declares no security schemes; AGNTCY across its SLIM, OASF, Directory, and identity specifications.}
  \label{tab:threat-surface}
  \footnotesize
  \begin{tabular}{lcccccccc}
    \toprule
    Threat class & A2A & MCP & ACP & AGNTCY & ANP & Coral & LOKA & ACNBP \\
    \midrule
    Prompt injection (protocol boundary) & N & N & N & P & N & N & P & \textbf{A} \\
    Agent impersonation & P & P & N & \textbf{A} & \textbf{A} & \textbf{A} & \textbf{A} & \textbf{A} \\
    Message replay & P & P & N & \textbf{A} & \textbf{A} & N & \textbf{A} & \textbf{A} \\
    Capability escalation & N & N & N & N & N & P & P & \textbf{A} \\
    Delegation token theft & N & N & N & N & P & P & P & \textbf{A} \\
    Registry poisoning & N & N & N & \textbf{A} & \textbf{A} & N & \textbf{A} & \textbf{A} \\
    AgentCard supply-chain & N & n/a & N & \textbf{A} & n/a & n/a & N & \textbf{A} \\
    Message tampering / MITM & \textbf{A} & \textbf{A} & \textbf{A} & \textbf{A} & \textbf{A} & P & \textbf{A} & \textbf{A} \\
    Side-channel privacy leak & N & N & N & N & N & N & P & P \\
    Denial of service & P & N & N & P & N & N & N & \textbf{A} \\
    Cross-agent privilege confusion & N & n/a & N & \textbf{A} & N & \textbf{A} & \textbf{A} & P \\
    Audit-log tampering & N & N & N & P & \textbf{A} & \textbf{A} & \textbf{A} & \textbf{A} \\
    \bottomrule
  \end{tabular}
\end{table}


The matrix renders three structural observations, and a
fourth that only appears once the twelve classes are
weighted by how much each one costs.

\paragraph{The widely-deployed protocols cluster at one
addressed class each.}
A2A, MCP, and ACP each address exactly one of the twelve
classes, and in all three cases it is message tampering, via
transport encryption. A2A adds three partial mitigations and
MCP two, so the specifications are not silent on security;
they are non-normative about it. The reason is consistent
across the cluster: each protocol
specifies its identity and transport-encryption layers,
the identity layer without binding it cryptographically,
then relies on deployment-level mitigations
(WAFs, OAuth scopes, application-layer audit logging) for
the rest. AGNTCY is the counterexample that shows this is a
choice rather than a constraint of industry sponsorship: it
is equally industry-led and addresses six classes, because
MLS~\cite{rfc9420}, Sigstore signing~\cite{sigstore2022}, and
content addressing are specified
normatively rather than recommended. The literature
converges on the assessment that
the cluster's posture is the principal residual risk in
production: A2ASecBench~\cite{a2asecbench} reports
exploitable variants of nine of the twelve classes against
deployed A2A endpoints in its 2026 evaluation; the A2A
weakness analysis~\cite{a2a_weaknesses} catalogues eight
named threats against the same spec.

\paragraph{Decentralised-identity protocols invert the
profile.}
LOKA (6 addressed), ANP (5 addressed), and Coral (3
addressed) trade spec coverage of operational
threats for coverage of identity, authentication, and
audit-log threats via DIDs, Verifiable Credentials, and (in
Coral's and LOKA's cases) blockchain-anchored audit
trails. The cluster's weakness is the inverse: prompt
injection, capability escalation, delegation token theft,
and denial of service are unaddressed or only partial
throughout. Coral is the weakest of the three on the count
because its cryptographic guarantees attach to the payment
path rather than the message path. The
DIDs-and-VCs-for-agents paper~\cite{did_vc_agents} treats
this trade-off explicitly, framing it as the consequence of
specifying identity rigorously but leaving execution-time
threats to higher layers.

\paragraph{ACNBP is a positive outlier.}
ACNBP's ten of twelve addressed classes, with the remaining
two partial rather than absent (side-channel privacy via
ring signatures, and cross-agent privilege confusion via
agent quarantine), is the highest spec-level threat
coverage in the surveyed set, achieved by the explicit
MAESTRO seven-layer threat-modelling pass within the
specification (\S\ref{sec:emerging:acnbp})~\cite{spec_acnbp}.
ACNBP's deployment maturity is the lowest in the cluster:
the protocol exists at arXiv-draft stage with no
formal SDO. The design is still instructive. It shows what
specification-level threat coverage looks like when it is
written as part of the protocol rather than retrofitted
afterwards. \S\ref{sec:security:open} returns to ACNBP as
the model against which the industry-led protocols' threat
coverage gap should be benchmarked.

\paragraph{Severity-weighted coverage separates protocols
the count places together.}
The three observations above all read the matrix by
counting. Counting is the wrong instrument for the question
a deployer actually asks, which is how much risk is left
over. Table~\ref{tab:severity-coverage} reports the
severity-weighted view defined in \S\ref{sec:rubric:axis8}.
Each class carries a weight of 3, 2, or 1 by band. The
table reports addressed weight over applicable weight, and
the last column reports the unaddressed weight in the
critical and high bands alone.

% Generated by scripts/gen_table6.py: do NOT hand-edit.
% Source: data/threat-coverage.yaml (severity block).

\begin{table}[t]
  \centering
  \caption{Severity-adjusted threat coverage, the companion view to the Axis 8 count of Table~\ref{tab:threat-surface}. Bands and integer weights are declared in the companion repository, not computed here: 5 critical classes at weight 3, 5 high at weight 2, 2 moderate at weight 1, for 27 points across the twelve. Band columns read \emph{addressed / applicable}, where \texttt{na} classes are excluded from both. SWC is severity-weighted coverage: addressed weight over applicable weight. Residual is the unaddressed weight in the critical and high bands only, counting \texttt{partial} as unaddressed exactly as Axis 8 does. Higher SWC is better; lower residual is better.}
  \label{tab:severity-coverage}
  \footnotesize
  \begin{tabular}{lcccccr}
    \toprule
    Protocol & Axis 8 & Crit. & High & Mod. & SWC & Resid. \\
             & (of 12) & (w.\ 3) & (w.\ 2) & (w.\ 1) & & \\
    \midrule
    A2A & 1 & 0/5 & 1/5 & 0/2 & 7\% & 23 \\
    MCP & 1 & 0/5 & 1/3 & 0/2 & 9\% & 19 \\
    ACP & 1 & 0/5 & 1/5 & 0/2 & 7\% & 23 \\
    AGNTCY & 6 & 2/5 & 4/5 & 0/2 & 52\% & 11 \\
    ANP & 5 & 2/5 & 3/4 & 0/2 & 48\% & 11 \\
    Coral & 3 & 1/5 & 2/4 & 0/2 & 28\% & 16 \\
    LOKA & 6 & 2/5 & 4/5 & 0/2 & 52\% & 11 \\
    ACNBP & 10 & 5/5 & 4/5 & 1/2 & 89\% & 2 \\
    \bottomrule
  \end{tabular}
\end{table}


Three things follow that the count does not show. First,
the one-addressed cluster is worse off than its count
suggests. A2A, MCP, and ACP each address one class, and in
every case it is message tampering. That class sits in the
high band, not the critical one. All three therefore
carry zero of five critical classes. The two industry
protocols with the widest deployment leave 23 and 19 points
of critical-and-high weight unaddressed at the
specification level. Second, MCP's slightly higher
percentage is an artefact of scope, not of strength. Two
classes are \texttt{na} for MCP, so its denominator is
smaller, and the residual column is the fairer comparison.
Third, the gap between ACNBP and everything else widens
rather than narrows. ACNBP addresses all five critical
classes, so its residual is 2 against a field where the
next best is 11.

The parity that the count suggests between AGNTCY and LOKA
survives the reweighting. It is worth saying why, since
a weighting scheme that flattened every distinction would
be doing no work. Both address two of five critical classes
and four of five high, from different directions. AGNTCY
gets there through MLS, Sigstore signing, and content
addressing. LOKA gets there through DIDs, post-quantum
signatures, and a blockchain-anchored trail. The equal weighted score records
a real equivalence in residual risk, reached by
non-overlapping means.

The severity view does not rank protocols and is not
combined with any other axis. It answers one question the
count cannot, and it inherits the count's limits. It scores
specifications, not deployments. A specification that
addresses a critical class badly scores the same as one
that addresses it well.

\subsection{Identity and authentication}
\label{sec:security:identity}

Identity is the foundation on which every other security
property in inter-agent protocols rests, and the dimension
on which the contemporary protocols differ most. The rubric
matrix already records the three-way split
(\S\ref{sec:comparison:matrix}): industry-led protocols at
OAuth-bearer-token identity (Axis 2 score 1), ACNBP at
hardened-bearer-tokens-plus-PKI (Axis 2 score 2), and the
DID-anchored protocols at decentralised verifiable
identity (Axis 2 score 3).

\paragraph{The OAuth-bearer-token failure mode.}
A2A, MCP, and ACP all sit at OAuth-2.x identity. The
specification mandates bearer tokens issued by a third-party
authorisation server, with mTLS and JWT-with-signing as
optional hardenings. The empirical security
literature~\cite{a2asecbench, a2a_weaknesses} reports that
the optional hardenings are rarely used in the surveyed
deployments, and that bearer-token theft via downstream
log capture, browser cache leakage, or compromised
intermediaries is the most-frequent attack against deployed
A2A and MCP endpoints. The remediation under discussion in
the AAIF working group~\cite{lf_aaif} is to elevate the
mTLS-plus-message-signing pattern from optional to required
for production deployments, which would move A2A and MCP to
the rubric's Axis 2 level 2 in a future version.

\paragraph{The DID-anchored alternative.}
ANP, Coral, and LOKA all use W3C DIDs paired with Verifiable
Credentials (or, in Coral's case, Solana wallet signatures
over the DID). The advantage at the identity layer is
direct: a DID document carries the public keys necessary to
verify a challenge-response from the agent, and no
third-party authorisation server can be compromised to
issue tokens for an agent without the agent's
private-key cooperation. The
literature~\cite{did_vc_agents, spec_anp, spec_loka} reports
no published bypass of the
DID-plus-VC-challenge-response pattern; the residual risks
move down to private-key compromise (which is outside the
protocol scope) and to capability-layer attacks
(\S\ref{sec:security:authz}).

\paragraph{The PKI-hardened middle ground.}
ACNBP's X.509-plus-OCSP-plus-ECC envelope structure
\cite{spec_acnbp} sits between the OAuth and DID camps. The
operational advantage is compatibility with existing
enterprise PKI deployments; the operational disadvantage is
the cost of running an OCSP infrastructure that scales to
agent-population sizes. The literature is too thin to
report on real-world deployments of ACNBP's identity layer
because no such deployments exist at the cutoff; the
analysis here rests on the specification's own threat
model in \S VII of the ACNBP paper.

\paragraph{The protocol-level remaining problem.}
None of the contemporary protocols, including ACNBP,
addresses the \emph{identity rotation} problem. An agent
that needs to rotate its identifier (whether for compromise
recovery, key rotation, or organisational migration)
breaks any in-flight tasks that hold references to the
prior identifier. The DID specification supports rotation
in principle through DID-document
updates~\cite{w3c_did_core}, but no agent
protocol specifies how an in-flight Task's identity
mid-flight migrates from the old identifier to the new.
This is a candidate first item on the open-problems
agenda (\S\ref{sec:security:open}).

\subsection{Authorization propagation across agent chains}
\label{sec:security:authz}

The most underdeveloped area of the contemporary protocols
is the authorisation layer. Section~\ref{sec:rubric:axis3}
and the per-cell derivation in §\ref{sec:comparison:derivation}
already recorded the gap: A2A, MCP, and ACP at scope tokens
with no standardised delegation credential (level 1), and
ANP, Coral, LOKA, and ACNBP at scoped capabilities or a
two-party binding ceremony (level 2). Nothing reaches
cryptographically-bound capability chains. The level-3
cell is empty across the scored set, which makes this the
only axis where the strongest protocol and the weakest
differ in degree of partial coverage rather than in
whether the problem is addressed at all.

A distinction matters here and the survey keeps it
throughout. Not standardising a cryptographic delegation
chain and prescribing token forwarding are different
claims, and an earlier version of this survey ran them
together. We contacted Microsoft about that wording, and
this paragraph is the correction. A2A does the first and not
the second. It states
security requirements, advertises authentication schemes in
the AgentCard, and defines an in-task authorization flow
with an auth-required task state, but it leaves
authorization policy and the choice of delegation mechanism
to the surrounding identity and security
architecture~\cite{spec_a2a}. The level-1 score records the
omission. The failure mode below is a property of what
deployments then build, and the evidence for it is
deployment evidence.

The exclusion is also explicit rather than tacit. A2A's
\S7.6.4 states in normative language that the protocol does
not define the scope, representation, validity, or revocation
semantics of an authorization credential, and forbids reading
the auth-required transition as authorization for any
operation (\S\ref{sec:industry:a2a})~\cite{spec_a2a}. The
protocol supplies authorization knobs so that many
authorization solutions can sit above it. We contacted Cisco
as well, which pointed us to this clause and supplied that
reading. That is a design
choice, and the survey does not score it as a defect of
intent. It does score the consequence: a receiver holding
only protocol-level evidence still cannot verify a
delegation chain. A deliberate exclusion and an accidental
omission leave the receiver in the same position, which is
what the axis measures, so the survey keeps the two apart in
prose and does not distinguish them in the number.

\paragraph{The token-forwarding failure mode.}
A failure mode documented for multi-agent deployments is
OAuth token forwarding: orchestrator A
receives an OAuth token from the user; A invokes remote
agent B and passes that same token onward; B may invoke C
and pass it onward again. The pattern is invisible to the
authorisation server, which sees only that the token is
being used; the chain of agents through which the token
flowed is not recoverable from the audit log. The threat
class is documented in the \emph{Authorization Propagation
in Multi-Agent AI Systems} analysis~\cite{authz_propagation}
and in the Agentic JWT design~\cite{agentic_jwt}, both of
which propose successor patterns.

\paragraph{The capability-binding alternative, and where it
stops.}
ACNBP's Capability Binding ceremony (Definition 10 of the
spec~\cite{spec_acnbp}) replaces token forwarding for a
single hop with a mutually-signed commitment between the
calling and called agent. The caller signs a request
describing the specific capability to be invoked and the
bound parameters. The callee counter-signs the binding, and
the resulting artifact is the authorisation token for that
specific call. The binding is cryptographically verifiable
without consulting an external authorisation server, which
is what token forwarding cannot offer.

The ceremony does not, however, compose into a chain, and
the survey previously read it as though it did. ACNBP
specifies the binding between one requester and one
provider. It does not specify how a provider that becomes
a requester carries forward the authority it received, nor
how a third agent verifies the first grant. A lead author
confirmed in correspondence that the protocol is a
two-party negotiation-and-binding protocol and that
multi-agent delegation lies outside its scope. So the
ceremony is a better one-hop primitive than a forwarded
bearer token, and it is not yet an answer to the
propagation problem. Composing it is the open work, and it
is the reason the level-3 cell of Axis 3 is empty across
every protocol surveyed.

\paragraph{The Agentic JWT proposal.}
The Agentic JWT design~\cite{agentic_jwt} adapts the
JWT-based delegation pattern to multi-agent contexts. An
agent receives a JWT scoped to a specific intent class and
audience; when delegating to a downstream agent, the
upstream agent issues a derived JWT bound to the original
JWT's identifier through a cryptographic chain. Each
downstream agent can verify the chain without consulting the
original authorisation server. The proposal is not yet
adopted into any of the surveyed protocols; the AAIF
working group is reportedly studying it as a candidate
extension for the next A2A revision.

\paragraph{The cross-tenant boundary.}
A separate but related problem is cross-tenant
authorisation. When a single A2A agent serves multiple
principals from multiple tenants, the protocol does not
specify how the agent's per-call authorisation state
isolates one tenant's mandates from another's. The SAGA
governance architecture~\cite{saga2026} proposes a
tenant-policy layer that sits above the protocol and
enforces the isolation, but the residual risk is that the
protocol itself does not require the layer to exist.
\S\ref{sec:governance} returns to the question as an
institutional matter.

\subsection{Prompt injection as a protocol-level attack}
\label{sec:security:injection}

Prompt injection is a class of attack that did not exist
in the FIPA-ACL era (\S\ref{sec:background:fipa-kqml}): a
deterministic message parser cannot be ``confused'' by
unexpected natural-language content because it ignores
content outside its grammar. With LLM-based agents, the
content of a message is interpreted by a probabilistic
reasoner whose behaviour is sensitive to the literal text
in ways that the protocol specification cannot enumerate.
Prompt injection is therefore a protocol-level attack in
the precise sense that no application-level filter can
fully prevent it; the protocol must include the
defensive structure as a normative requirement.

\paragraph{The InjecAgent empirical baseline.}
Zhan et al.'s InjecAgent benchmark~\cite{zhan2024injecagent}
established that indirect prompt injection through tool
results is a measurable, reproducible threat across every
tool-binding deployment they evaluated. The
deployment-level mitigation (LLM-based filtering of
tool-output strings) is reported to remove the most
trivial attacks but leaves a long tail of more subtle
embeddings. The protocol-level mitigation that would
address the class is structured tool-output: typed
content blocks with content-negotiation guarantees
between the tool and the orchestrator. MCP's level 2 on
the rubric's Axis 7 is the closest the cutoff specs come
to this discipline; full level 3 (typed parts with
inline-or-reference duality and capability-negotiated
content) would substantially raise the bar.

\paragraph{The cross-agent injection vector.}
A second injection vector is cross-agent: an attacker who
controls remote agent B can craft a response to
orchestrator A that contains adversarial natural-language
content. The attack is conceptually the same as
tool-output injection but the mitigation is harder
because the response is opaque task output rather than
a structured tool result. A2A's DataPart content model
permits structured-data responses
but does not require them; the current convention is to
emit free-text. The
``From Prompt Injections to Protocol Exploits''
analysis~\cite{protocol_exploits} reports cross-agent
injection chains against deployed A2A endpoints and
recommends that A2A elevate DataPart from optional to
required for inter-agent message bodies, leaving TextPart
for the user-facing summary only.

\paragraph{The DECP partial mitigation.}
LOKA's Dynamic Ethics and Constraint
Protocol~\cite{spec_loka} is the only surveyed protocol
that includes prompt-injection filtering as part of the
specification, scored partial on Axis 8 because the
mechanism is described conceptually rather than
operationally. The mechanism applies a contextual-ethics
filter at the message-receive boundary, refusing to act
on messages whose intent fails ethical-policy validation.
The specification's threat-model treatment is more
explicit than the operational guidance, which is the
reason for the partial rather than addressed marker.

\paragraph{The MAESTRO foundation-model-layer treatment.}
ACNBP's MAESTRO threat analysis (\S VII.I of the
spec~\cite{spec_acnbp}) treats the foundation-model layer
as one of seven analysis layers and prescribes
spec-level mitigations including: a content-quarantine
boundary for untrusted message content, a
capability-scoped sandbox for tool outputs consumed by
the orchestrator, and a mandatory structured-output
format for inter-agent task results that prevents
free-text content from entering the orchestrator's
reasoning surface. The cost is operational complexity;
the benefit is the only spec-level treatment of prompt
injection across the surveyed set.

\subsection{Multi-agent privacy leakage}
\label{sec:security:privacy}

Privacy in multi-agent systems is structurally different
from privacy in single-agent systems because the message
boundary is no longer the trust boundary. When agent A
delegates a task to agent B, A's input to B may carry
context information that A's principal did not authorise
B to receive. The AgentLeak full-stack
benchmark~\cite{agentleak} establishes that this leakage
is a measurable, reproducible threat across every
multi-agent deployment they evaluated, and that the
typical leakage vector is conversation-history forwarding
in task-handoff payloads.

For evaluation, two dimensions must remain distinct.
\emph{Lifecycle stages} describe when data is handled
(input, candidate-pool construction, selection, execution,
and output); \emph{channels} describe where AgentLeak
observes it. In the AgentLeak trace model,
\texttt{user\_input} and \texttt{tool\_response} are source
channels. The core disclosure channels for forwarding
experiments are \texttt{inter\_agent\_message},
\texttt{shared\_memory}, \texttt{tool\_call}, and
\texttt{final\_output}; \texttt{log} and
\texttt{generated\_file} are additional operational
disclosure channels when instrumented. A selection record
is therefore joined to later channel events by candidate
lineage; it is not relabelled as a disclosure channel.

\paragraph{The conversation-history vector.}
Across the coordinator-worker configurations evaluated by
AgentLeak, forwarding prior conversation history to a
remote agent is a recurrent leakage vector. The convenience
advantage is direct: the remote agent does not need to ask
clarifying questions. The privacy cost is that the
history may include content from prior turns that the
principal would not have shared with the remote agent if
asked directly. The protocol specifications do not
require the orchestrator to filter history, and the
AgentLeak benchmark~\cite{agentleak} reports that
\textbf{inter-agent messages leak personally-identifiable
content at 68.8\%}, compared to 27.2\% on the
multi-agent output channel and 43.2\% on the single-agent
output channel. Total-system exposure across the four
evaluated domains (healthcare, finance, legal, corporate)
reaches 68.9\%. The asymmetry is the central empirical
finding: inter-agent communication is the dominant
leakage vector, and output-only audits miss roughly
41.7\% of violations~\cite{agentleak}.

\paragraph{The contextual-integrity model.}
The literature converges on Nissenbaum's contextual
integrity~\cite{nissenbaum2009privacy} as the conceptual
framework: information should flow only according to
norms appropriate to the original context. Kong et
al.~\cite{kong2025communication} is where the framework
enters the inter-agent protocol literature. The protocol
implication is that the orchestrator's
forwarding decision should be context-sensitive.
Forwarding a billing query to a finance agent is
appropriate; forwarding it to a calendar agent is not.
None of the contemporary protocols specifies a
contextual-integrity primitive; the open question is
whether the primitive belongs at the protocol layer or
at the orchestrator's policy layer. The survey's own
candidate for the protocol-layer answer is the ACSP
context selection envelope of \S\ref{sec:proposals:acsp}
(Contribution C7), which specifies a wire-level audit log
on the sender's selection step that the receiver can
verify; we hold the orchestrator-versus-protocol dispatch
open in \S\ref{sec:openproblems:regulation} (P5).

An implementation study of this primitive must report the
hard boundary decision separately from graded risk and
task utility. Per-channel incidence, entity type, severity,
and repeated propagation are diagnostics of observed
disclosure; none can substitute for the allow/deny policy
that governs whether an item may cross a boundary. An
unobserved channel is missing evidence, not a clean result.

\paragraph{Side-channel leakage.}
Beyond the content-bearing channels measured by AgentLeak,
a separate open privacy threat is side-channel leakage through
metadata: task lifecycle timestamps, retry counts, error
messages, and AgentCard metadata can together reveal
information about the principal's actions even when the
message content is encrypted. This survey's specification
review found no explicit treatment of that metadata-flow
threat in the surveyed protocols, even
in the strongest cases (LOKA's contextual-ethics
policies cover content but not metadata; ACNBP's
ring signatures address sender privacy but not action
privacy).

\subsection{Confidential computing and TEE-based hardening}
\label{sec:security:tee}

The most recent direction in the literature is the
proposal that orchestrators and remote agents should run
inside Trusted Execution Environments (TEEs), such as
Intel SGX, AMD SEV-SNP, Apple Secure Enclave, or AWS
Nitro Enclaves. The protocol would carry remote
attestation of the TEE's measurement as part of agent
identity. The
\emph{When Agents Handle Secrets} survey~\cite{confidential_computing}
catalogues the proposal space and identifies three
distinct claims it can make: (a) the orchestrator's
reasoning surface is unobservable to the host,
(b) the orchestrator's key material is unextractable, and
(c) the orchestrator's policy
configuration is tamper-evident.

\paragraph{What TEEs solve.}
Within the threat model of \S\ref{sec:security:model},
TEEs are well-matched to the
orchestrator-implementation-compromise threat: an
adversary who controls the host on which the
orchestrator runs cannot inspect the orchestrator's
prompts, intermediate reasoning, or tool-call payloads.
The remote-attestation protocol allows downstream agents
to verify that they are communicating with a
TEE-resident counterpart rather than a host-resident
impostor.

\paragraph{What TEEs do not solve.}
TEEs do not address the principal-orchestrator threat
boundary (the principal still has to trust the
orchestrator's prompt to it), the
orchestrator-remote-agent threat boundary (the remote
agent is still itself an untrusted entity from the
orchestrator's perspective), or the
foundation-model layer threats. The literature is in
agreement~\cite{confidential_computing, trism} that TEEs
are a useful component in a defence-in-depth architecture
but cannot stand alone.

\paragraph{The TRiSM framework.}
The TRiSM framework for agentic AI~\cite{trism} provides
the architectural integration: identity (DIDs), authorisation
(capability chains), runtime trust (TEEs and attestation),
and observability (audit logs) compose into a single
governance posture. The framework is descriptive rather
than prescriptive: it does not propose a new protocol.
It still usefully aligns the dispersed
literature around a four-pillar model.

\paragraph{The SAGA architecture.}
The SAGA paper~\cite{saga2026} (NDSS 2026) is the most
concrete proposal for a TEE-anchored governance layer
above the contemporary protocols. SAGA proposes that
each tenant's policy be enforced by a TEE-resident
gateway that validates every inter-agent message against
the tenant's policy before passing the message to the
underlying protocol. The cost is a per-message latency
hit; the benefit is that tenant policy is enforced
cryptographically rather than relying on the
underlying protocol to comply. SAGA is the most cited
recent proposal for the residual cross-tenant problem
identified at the end of \S\ref{sec:security:authz}.

\subsection{Open security problems}
\label{sec:security:open}

The security literature, taken together, identifies six
unresolved problems that no contemporary protocol fully
addresses. We list them with their priority ranking from
the literature; \S\ref{sec:openproblems} returns to half
of them as the highest items on the survey's
open-problems agenda.

\paragraph{O1: Cryptographic delegation chains.}
The single most-cited gap, and the one no surveyed
protocol closes. No protocol standardises a delegation
credential, so token forwarding is the de-facto pattern
deployments adopt, and it is recoverable only by a
successor pattern that carries verifiable
chain-of-custody information. The Agentic
JWT~\cite{agentic_jwt} and authorisation-propagation
\cite{authz_propagation} proposals are the leading
candidates. ACNBP's Capability Binding ceremony
\cite{spec_acnbp} is the most rigorous spec-level
treatment of a single hop, and composing single hops into
a verifiable chain is exactly the part it leaves open.

\paragraph{O2: Identity rotation.}
None of the contemporary protocols specifies how an
agent's identity rotates without breaking
in-flight tasks. The problem is operationally critical
for compromise recovery and key rotation. No published
proposal at the cutoff fully addresses it.

\paragraph{O3: Prompt-injection mitigation as a
spec-level requirement.}
Only ACNBP's MAESTRO treatment~\cite{spec_acnbp} and
LOKA's DECP~\cite{spec_loka} include prompt-injection
mitigation as part of the specification. The
industry-led protocols leave it as an
application-level concern. The next protocol revision
in the AAIF working group~\cite{lf_aaif} should
include a normative content-quarantine boundary at the
message-receive layer.

\paragraph{O4: Cross-tenant policy enforcement.}
Multi-tenant agentic deployments require per-tenant
policy isolation that the contemporary protocols do not
provide. SAGA~\cite{saga2026} is the leading proposal
for a TEE-anchored gateway architecture; the open
question is whether the gateway is a protocol concern
or an infrastructure concern.

\paragraph{O5: Contextual integrity for
multi-agent forwarding.}
The conversation-history forwarding pattern is the
dominant privacy leak in production multi-agent
deployments~\cite{agentleak}. No contemporary protocol
specifies a contextual-integrity primitive that would
let the orchestrator make forwarding decisions
according to context-appropriate norms.

\paragraph{O6: Side-channel metadata privacy.}
Even when content is encrypted, the metadata associated
with inter-agent messages (lifecycle timestamps, error
patterns, retry counts) leaks information about the
principal's actions. ACNBP's ring signatures address
sender privacy partially; no protocol addresses
action privacy. The literature is thin on this front,
which is itself an indicator of how early the field
is in treating side-channel threats.

These six problems together comprise what we treat in
\S\ref{sec:openproblems} as the security agenda for the
next standardisation cycle.
